% Capitulo: Modelagem para ciência de dados
\chapter{Modelagem para ciência de dados}
\label{cap:modelagem-cd}

\rascunho{
  \textbf{Ideia do capítulo.} Análises preditivas e prescritivas pedem uma modelagem
  diferente da descritiva. O esquema estrela é ótimo para somar e filtrar, mas um modelo
  preditivo precisa de uma linha por entidade e por momento, com o cuidado de só usar o
  que se sabia naquele momento. E a prescritiva ainda precisa guardar as decisões
  tomadas.
  \begin{itemize}
    \item Os três níveis (descritivo, preditivo, prescritivo) como fio do capítulo
      \parencite{davenport2007}. \pendente{conferir se uso essa referência ou outra mais recente para os três níveis}
    \item Posso ligar com a minha experiência de pesquisa, que segue os três níveis.
      \pendente{decidir se cito o mestrado}
    \item Referência principal para a parte técnica: \textcite{huyen2022}.
  \end{itemize}
}

% ---------------------------------------------------------------------------
\section{Descritiva, preditiva e prescritiva}
\label{sec:tres-niveis}

\rascunho{
  \begin{itemize}
    \item Descritiva: o que aconteceu. Preditiva: o que vai acontecer. Prescritiva: o
      que fazer a respeito.
    \item Cada nível depende do anterior: não se prevê bem o que não se descreve bem.
    \item Um exemplo que atravessa os três níveis. \pendente{escolher o domínio: agricultura, finanças ou logística}
  \end{itemize}
}

\pendente{escrever}

% ---------------------------------------------------------------------------
\section{A tabela de variáveis}
\label{sec:tabela-features}

\rascunho{
  \begin{itemize}
    \item O grão muda: em vez de \enquote{uma linha por venda}, \enquote{uma linha por
      cliente por mês} ou \enquote{uma linha por veículo por dia}, com as variáveis
      (\emph{features}) calculadas e o alvo que se quer prever.
    \item Tabelas largas, com muitas colunas, que o modelo dimensional evitaria.
    \item A camada prata costuma ser a fonte melhor que o ouro, porque o ouro já vem
      agregado para o BI e perde detalhe.
  \end{itemize}
}

\pendente{escrever}

% ---------------------------------------------------------------------------
\section{O futuro não pode vazar para o passado}
\label{sec:vazamento}

\rascunho{
  \begin{itemize}
    \item O erro mais comum e mais caro: calcular uma variável com dado que só existiu
      depois do momento da previsão (vazamento, \emph{data leakage}) \parencite{huyen2022}.
      O modelo fica ótimo no teste e inútil em produção.
    \item Correção no tempo (\emph{point-in-time}): cada variável é calculada só com o
      que se sabia na data de referência. A consulta abaixo mostra a ideia.
    \item O histórico das dimensões (as dimensões que mudam devagar, de Kimball) vira
      requisito: é preciso saber como o cadastro estava naquele dia, não como está hoje.
  \end{itemize}
}

\pendente{escrever}

\begin{sql}
-- Variáveis de cada cliente calculadas só com o passado de cada data de referência
SELECT
    r.id_cliente,
    r.data_referencia,
    COUNT(v.id_venda)                AS compras_90d,
    COALESCE(SUM(v.valor_total), 0)  AS valor_90d,
    r.cancelou_no_mes_seguinte       AS alvo
FROM referencias_clientes AS r
LEFT JOIN vendas AS v
       ON v.id_cliente = r.id_cliente
      AND v.data_venda <  r.data_referencia                       -- nada do futuro
      AND v.data_venda >= r.data_referencia - INTERVAL '90' DAY
GROUP BY r.id_cliente, r.data_referencia, r.cancelou_no_mes_seguinte;
\end{sql}

% ---------------------------------------------------------------------------
\section{Reuso de variáveis e dívida técnica}
\label{sec:feature-store}

\rascunho{
  \begin{itemize}
    \item Quando cada cientista calcula as próprias variáveis, o mesmo conceito ganha
      várias definições, como acontecia com as métricas antes do data warehouse.
    \item Repositório de variáveis (\emph{feature store}): definir uma vez, usar no
      treino e em produção com o mesmo código \parencite{huyen2022}.
    \item A maior parte do código de um sistema de aprendizado de máquina não é o
      modelo, é dado e infraestrutura em volta dele \parencite{sculley2015}.
  \end{itemize}
}

\pendente{escrever}

% ---------------------------------------------------------------------------
\section{Da previsão à decisão}
\label{sec:prescritiva}

\rascunho{
  \begin{itemize}
    \item A prescritiva junta a previsão com o que o negócio pode fazer: custos,
      capacidades, restrições, regras. Muitas vezes é um problema de otimização.
    \item O dado novo que ela exige: o registro da decisão recomendada, da decisão
      tomada e do resultado. Sem isso não se sabe se a recomendação funciona.
    \item Modelar cenários (\enquote{e se}) como dados, não como planilhas soltas.
    \item Ligação com o capítulo~\ref{cap:negocio}: as restrições certas vêm de quem
      conhece o negócio.
  \end{itemize}
}

\pendente{escrever}

% ---------------------------------------------------------------------------
\begin{leituras}
  \leitura{huyen2022}{A visão mais completa e prática de como dados, variáveis e
    modelos se organizam num sistema de aprendizado de máquina em produção.}
  \leitura{sculley2015}{Artigo curto que mostra por que o modelo é a menor parte do
    problema e onde a dívida técnica se esconde.}
  \leitura{davenport2007}{Para a visão de negócio dos níveis de análise e de como as
    empresas competem usando dados.}
\end{leituras}
